% 图契约：本题的源、边界和推导逐项对应；箭头符号按图注解释。
\documentclass[tikz,border=7pt]{standalone}
\usepackage{fontspec,xeCJK,amsmath,amssymb}
\IfFontExistsTF{Times New Roman}{\setmainfont{Times New Roman}}{\setmainfont{TeX Gyre Termes}}
\IfFontExistsTF{Heiti SC}{\setCJKmainfont[AutoFakeBold=1.5]{Heiti SC}}{\setCJKmainfont[AutoFakeBold=1.5]{FandolHei-Regular}}
\usetikzlibrary{arrows.meta,bending,calc,patterns,positioning}
\definecolor{ink}{HTML}{243745}
\definecolor{blue}{HTML}{275DAB}
\definecolor{green}{HTML}{177C63}
\definecolor{red}{HTML}{BD4935}
\definecolor{orange}{HTML}{AA6C16}
\definecolor{purple}{HTML}{715393}
\tikzset{every node/.style={font=\fontsize{9.3}{12.5}\selectfont,text=ink},
axis/.style={-{Stealth[length=4pt]},line width=.65pt,draw=ink!65},
vec/.style={-{Stealth[length=5pt]},line width=1pt,draw=blue},
short/.style={-{Stealth[length=3pt]},line width=.8pt},
guide/.style={dashed,draw=ink!40,line width=.5pt},
card/.style={rounded corners=2pt,draw=blue!65,fill=blue!4,align=center,text width=6.5cm,inner xsep=9pt,inner ysep=7pt},
title/.style={font=\fontsize{10.5}{14}\selectfont\bfseries,text width=15cm,align=center},
label/.style={font=\fontsize{8.4}{11}\selectfont,align=center},
note/.style={label,text width=6.2cm}}
\begin{document}
\begin{tikzpicture}[x=1cm,y=1cm]
\path[use as bounding box] (-8,-3.5) rectangle(8,3.5);
\node[title] at (0,3.15) {M2-022\quad 非均匀场使两端电力不能抵消};
\begin{scope}[shift={(-4,0)}]
\draw[axis] (-2.3,0)--(2.4,0) node[right] {$x$};
\fill[blue] (-0.95,0) circle(2.3pt) node[above=3pt] {$-q$};
\fill[red] (0.95,0) circle(2.3pt) node[above=3pt] {$+q$};
\draw[vec,green] (-.7,-.8)--(.7,-.8) node[midway,below] {$\mathbf d$};
\draw[vec](-1.8,1.3)--(-1,1.3);\draw[vec](.2,1.3)--(1.8,1.3);\node[label]at(0,2){$E$ 向右增强（$\gamma>0$）};\node[note] at (0,-2.85) {图示 $p>0,\gamma>0$；通式按实际系数取符号。};
\end{scope}
\node[card,draw=blue!65,fill=blue!4] at (3.9,1.8) {$\mathbf F=(\mathbf p\cdot\nabla)\mathbf E$};
\node[card,draw=green!65,fill=green!4] at (3.9,0) {$\mathbf F=p\gamma\mathbf e_x$};
\node[card,draw=orange!65,fill=orange!4] at (3.9,-1.8) {正端的右向力较大；\\负端的左向力较小。};
\end{tikzpicture}\end{document}